% !TeX program = pdflatex
%==============================================================================
%  Verificação de Reposição (VR) --- prova de segunda chamada, aulas 01 a 10
%
%  O gabarito sai de "Verificacao de reposicao - gabarito.tex", que apenas liga
%  a opção [gabarito] e inclui este arquivo --- fonte única, dois PDFs.
%
%  QUATRO QUESTÕES, DE 3,0 PONTOS: duas sorteadas da "Lista de revisao 01-06" e
%  duas da "Lista de revisao 07-10", uma aula por questão dentro de cada lista.
%  Ficaram fora do sorteio os exercícios que já tinham caído na Avaliação Teórica
%  01 (Ex. 5, 7 e 11 da 01-06) e na 02 (Ex. 3, 6 e 9 da 07-10). A VR e a
%  Verificação Suplementar (VS) foram sorteadas juntas, a VR primeiro, e não têm
%  exercício em comum. O sorteio é reproduzível:
%
%      import numpy as np
%      rng = np.random.default_rng(20261116)                 # a data, aaaammdd
%      # exercícios que ainda não caíram em prova, por aula
%      lista_01_06 = {1: [1, 2], 2: [3, 4], 3: [6], 4: [8], 5: [9, 10]}
%      lista_07_10 = {7: [1, 2], 8: [4], 9: [5], 10: [7, 8]}
%
%      def sorteia(lista):
%          aulas = sorted(rng.choice(sorted(lista), size=2, replace=False).tolist())
%          escolhidos = []
%          for a in aulas:
%              e = int(rng.choice(lista[a]))
%              escolhidos.append(e)
%              lista[a].remove(e)
%              if not lista[a]:
%                  del lista[a]
%          return escolhidos
%
%      vr = sorteia(lista_01_06) + sorteia(lista_07_10)
%      vs = sorteia(lista_01_06) + sorteia(lista_07_10)
%      # -> vr = [1, 4, 2, 5] (aulas 01, 02, 07 e 09)
%      #    vs = [2, 3, 1, 7] (aulas 01, 02, 07 e 10)
%
%  Resultado, para a VR:
%      Aula 01 -> Ex. 1 da 01-06  ex-01-preditor-otimo (o melhor preditor)
%      Aula 02 -> Ex. 4 da 01-06  ex-04-oraculo (o que o MQO estima)
%      Aula 07 -> Ex. 2 da 07-10  ex-02-preco-ingenuo (o preço do ingênuo)
%      Aula 09 -> Ex. 5 da 07-10  ex-05-vetores-suporte (a margem máxima e os
%                                 vetores de suporte)
%
%  O sorteio rodou uma vez, em 16/11/2026, e vale. O algoritmo é o das
%  Avaliações Teóricas 01 e 02, com duas aulas por lista e a retirada de cada
%  exercício sorteado.
%
%  NÃO HÁ TEXTO DE EXERCÍCIO NESTE ARQUIVO. Os enunciados e as soluções vêm por
%  \input dos MESMOS arquivos que as listas de revisão usam, em exercicios/ e em
%  exercicios-07-10/. Os pontos por item escritos neles somam 4,0, o peso de uma
%  questão nas provas de três questões; aqui, o \pontositens antes de cada
%  \input dá os valores desta prova, que somam 3,0.
%
%  A PROVA NÃO REVELA DE QUE AULA VEM CADA QUESTÃO: o \questao não leva marcador
%  de aula, os títulos são os da lista (vazios onde a lista não tem título), e a
%  caixa de instruções traz a notação de regressão e a de classificação juntas.
%==============================================================================
\providecommand{\gabaritoopt}{}
\documentclass[11pt,a4paper]{article}
\usepackage[\gabaritoopt]{../recursos/latex/estilo-avaliacao}

\provatrue
\renewcommand{\tituloprova}[1]{Verificação de Reposição (#1)}

\begin{document}
\cabecalhoprova[{salvo indicação em contrário, $(\X,Y)$ é um vetor aleatório com
            $\X\in\R^p$, e $p$ é o número de covariáveis;
            $\Dados=\{(\X_i,Y_i)\}_{i=1}^n$ é a amostra de treino, i.i.d. Em
            regressão, $r(\x)=\E[Y\mid\X=\x]$ é a função de regressão, e a perda é a
            quadrática. Em classificação, $Y$ assume valores num conjunto finito
            $\mathcal{C}$ de classes; no caso binário, $\mathcal{C}=\{0,1\}$, a
            classe $Y=1$ é a positiva e $\eta(\x)=\Prob(Y=1\mid\X=\x)$, e a perda é
            a 0--1, com risco $\risco(g)=\Prob\big(Y\ne g(\X)\big)$.}]{VR}{01 a 10}{Regressão e Classificação}{2 horas}{12,0 pontos}

%==============================================================================
\questao{O melhor preditor}{3,0 pontos}
%==============================================================================
\pontositens{0,9;0,7;0,6;0,8}
% Aula 01 --- o preditor ótimo e o erro irredutível.
% Incluído por "Lista de revisao 01-06.tex" e, se sorteado, pela prova.
No início do curso, vimos uma afirmação forte: de todas as funções que você poderia usar
para prever $Y$ a partir de $\X$, a melhor é a função de regressão
$r(\x)=\E[Y\mid\X=\x]$. Não a melhor entre as lineares, nem a melhor entre as que
o seu método consegue produzir --- a melhor, ponto.

Uma afirmação dessas pede demonstração, e ela é mais curta do que parece. Todo o
trabalho está num truque que você vai reencontrar várias vezes no curso: somar e
subtrair a quantidade certa dentro de um quadrado. Ao longo do exercício, suponha
$\E[Y^2]<\infty$.
\begin{itens}
  \item \pts{1,2} Seja $g:\R^p\to\R$ uma função qualquer com
        $\E[g(\X)^2]<\infty$. Some e subtraia $r(\X)$ dentro do quadrado e mostre
        que
        \[
          \E\big[(Y-g(\X))^2\big]
            = \E\big[(Y-r(\X))^2\big] + \E\big[(r(\X)-g(\X))^2\big].
        \]
        O trabalho está em mostrar que o termo cruzado desaparece; condicione em
        $\X$ e pergunte-se o que $\E[Y-r(\X)\mid\X]$ vale.
  \item \pts{0,9} Olhe para a identidade que você acabou de provar e responda: por
        que ela mostra que $r$ é o minimizador? Em seguida, mostre que o valor
        mínimo é $\risco(r)=\E[\Var(Y\mid\X)]$.
  \item \pts{0,8} Esse mínimo tem nome: \textbf{erro irredutível}. O nome sugere
        que não há como baixá-lo, e é verdade --- mas a razão não é que os métodos
        atuais sejam ruins. Explique de onde vem essa impossibilidade, e diga o que
        teria de mudar \emph{no problema} para o número diminuir.
  \item \pts{1,1} Sabe-se que o melhor preditor \emph{constante} tem risco
        $\Var(Y)$. Use a identidade de (a) para escrever
        a diferença $\Var(Y)-\risco(r)$ como uma esperança, e interprete: que
        pergunta prática essa quantidade responde?
\end{itens}

\begin{solucao}
Como supõe o enunciado, $\E[Y^2]<\infty$: é o que garante que todas as esperanças
abaixo existam. Em particular $r(\X)$ também tem segundo momento
finito, porque $\E[r(\X)^2]=\E\big[(\E[Y\mid\X])^2\big]\le\E\big[\E[Y^2\mid\X]\big]=\E[Y^2]$
(desigualdade de Jensen, aplicada condicionalmente a $\X$).

\smallskip
\textbf{(a)} Some e subtraia $r(\X)$ e expanda o quadrado:
\[
  (Y-g(\X))^2 = \big(Y-r(\X)\big)^2
     + 2\big(Y-r(\X)\big)\big(r(\X)-g(\X)\big)
     + \big(r(\X)-g(\X)\big)^2 .
\]
As três parcelas são integráveis: as duas quadráticas porque $Y$, $r(\X)$ e $g(\X)$
têm segundo momento finito, e o termo cruzado pela desigualdade de Cauchy--Schwarz,
$\E\abs{UV}\le\sqrt{\E[U^2]\,\E[V^2]}$. Podemos, então, tomar esperança termo a
termo, e tudo se resume a mostrar que o termo cruzado tem esperança zero.

Condicione em $\X$ e use a lei da esperança total:
\[
\begin{aligned}
  \E\big[(Y-r(\X))(r(\X)-g(\X))\big]
  &= \E\Big[\,\E\big[(Y-r(\X))(r(\X)-g(\X))\,\big|\,\X\big]\Big]\\
  &= \E\Big[\big(r(\X)-g(\X)\big)\,
       \underbrace{\E\big[Y-r(\X)\,\big|\,\X\big]}_{=\;r(\X)-r(\X)\;=\;0}\Big]
  = 0 .
\end{aligned}
\]
O passo do meio é a regra "o que é função de $\X$ sai da esperança condicional em
$\X$": $r(\X)-g(\X)$ depende só de $\X$, e condicionar em $\X$ é tratá-lo como
conhecido. O zero vem da \emph{definição} de $r$: $\E[Y\mid\X]=r(\X)$.

Em palavras: o erro $Y-r(\X)$ que sobra depois de acertar a média condicional é
\textbf{ortogonal a toda função de $\X$} --- nenhuma informação contida em $\X$ ajuda a
prevê-lo. Guarde a frase: ela tem uma versão para o melhor preditor linear, cujo
resíduo é ortogonal a cada covariável.

\smallskip
\textbf{(b)} A identidade de (a) diz que
\[
  \risco(g) \;=\; \underbrace{\risco(r)}_{\text{não depende de } g}
   \;+\; \underbrace{\E\big[(r(\X)-g(\X))^2\big]}_{\ge\ 0} .
\]
Logo $\risco(g)\ge\risco(r)$ para toda $g$: $r$ é minimizadora. E é a única, a menos
de probabilidade zero: a igualdade vale se e só se $\E[(r(\X)-g(\X))^2]=0$, isto é,
se $g(\X)=r(\X)$ com probabilidade $1$. (Alterar $g$ num conjunto de valores de $\x$
que nunca ocorrem não muda risco nenhum.)

O valor mínimo sai de novo pela esperança total. Dado $\X$, $Y-r(\X)$ é o desvio de
$Y$ em torno da sua própria média condicional, e a esperança condicional do seu
quadrado é, por definição, a variância condicional:
\[
  \risco(r)=\E\big[(Y-\E[Y\mid\X])^2\big]
   =\E\Big[\,\E\big[(Y-\E[Y\mid\X])^2\;\big|\;\X\big]\Big]
   =\E\big[\Var(Y\mid\X)\big].
\]

\smallskip
\textbf{(c)} A impossibilidade é matemática, e não tecnológica. O item (b) mostrou que
\textbf{nenhuma} função de $\X$ tem risco menor que $\E[\Var(Y\mid\X)]$ --- nem mesmo
$r$, conhecida exatamente, com $n=\infty$ e erro de estimação zero. E essa quantidade
não menciona método nenhum: é uma propriedade da distribuição conjunta de $(\X,Y)$.
Ela mede quanto $Y$ ainda varia entre indivíduos com \emph{o mesmo} $\x$. Dois
pacientes com idade, peso e exames idênticos, mas desfechos diferentes, recebem
forçosamente a mesma previsão, porque o preditor recebeu a mesma entrada nos dois
casos.

Para baixar o número é preciso mudar o problema: \textbf{medir mais}. Se $Z$ é uma
covariável nova, a lei da variância total, aplicada condicionalmente a $\X$, dá
\[
  \E\big[\Var(Y\mid\X)\big]
  = \E\big[\Var(Y\mid\X,Z)\big]
    + \E\Big[\Var\big(\E[Y\mid\X,Z]\;\big|\;\X\big)\Big]
  \;\ge\; \E\big[\Var(Y\mid\X,Z)\big],
\]
com igualdade só quando $Z$ não muda a média condicional, $\E[Y\mid\X,Z]=\E[Y\mid\X]$.
Acrescentar informação nunca aumenta o erro irredutível, e o diminui sempre que a
informação for relevante. Isso vale na população; com $n$ finito, cada covariável
nova também custa variância de estimação, de modo que medir mais é
necessário, mas não suficiente, para prever melhor.

\smallskip
\textbf{(d)} Aplique a identidade de (a) com $g\equiv\E[Y]$. O lado esquerdo é o risco
do melhor preditor constante, $\Var(Y)$. E como $\E[r(\X)]=\E\big[\E[Y\mid\X]\big]=\E[Y]$,
pela esperança total, a última parcela é uma variância:
\[
  \Var(Y) = \risco(r) + \E\Big[\big(r(\X)-\E[r(\X)]\big)^2\Big]
  \qquad\Longrightarrow\qquad
  \boxed{\;\Var(Y)-\risco(r)=\Var\big(r(\X)\big)\;}
\]
É a lei da variância total, $\Var(Y)=\E[\Var(Y\mid\X)]+\Var(\E[Y\mid\X])$, obtida por
outro caminho.

A pergunta que isso responde é: \textbf{vale a pena usar as covariáveis?} O lado
esquerdo é o máximo que se pode ganhar trocando o preditor constante pelo melhor
preditor possível; o direito diz que esse ganho é exatamente o quanto a resposta
\emph{esperada} muda de um perfil de covariáveis para outro. Dividido por $\Var(Y)$,
vira a fração da variância de $Y$ que as covariáveis conseguem explicar no melhor dos
casos: um $R^2$ populacional, teto para o $R^2$ de qualquer modelo.

Os dois extremos ajudam a fixar. Se $\Var(r(\X))=0$, a função de regressão é
constante: as covariáveis não carregam informação sobre a média de $Y$, e prever pela
média é o melhor que existe --- não há modelo a construir. Se $\risco(r)=0$, $Y$ é
função de $\X$, e toda a sua variância é explicável.
\end{solucao}


\begin{criterio}
\textbf{(a)} 0,2 por somar e subtrair $r(\X)$ e expandir o quadrado; 0,5 pelo termo
cruzado: 0,2 por condicionar em $\X$ (esperança total), 0,2 por tirar
$r(\X)-g(\X)$ da esperança condicional e 0,1 por $\E[Y-r(\X)\mid\X]=0$, pela definição
de $r$; 0,2 pela identidade final. Afirmar que o termo cruzado se anula sem
condicionar em $\X$ vale 0,1 desses 0,5. A integrabilidade das parcelas não é exigida.

\textbf{(b)} 0,3 por ler na identidade que $\risco(g)=\risco(r)+\E[(r(\X)-g(\X))^2]$,
com a primeira parcela independente de $g$ e a segunda não negativa, e concluir
$\risco(g)\ge\risco(r)$; 0,4 pelo valor mínimo: 0,2 por condicionar em $\X$ e 0,2 por
reconhecer a variância condicional. A unicidade, a menos de probabilidade zero, não é
exigida.

\textbf{(c)} 0,3 pela origem da impossibilidade: o mínimo vale para qualquer função de
$\X$, mesmo para $r$ conhecida exatamente, e é uma propriedade da distribuição de
$(\X,Y)$, a variação de $Y$ entre indivíduos com o mesmo $\x$. 0,3 pelo que teria de
mudar no problema: medir covariáveis novas, que reduzem o erro irredutível quando mudam
a média condicional. Responder ``mais dados'', isto é, aumentar $n$, não pontua nesses
0,3, porque $n$ não aparece no erro irredutível.

\textbf{(d)} 0,3 por aplicar a identidade de (a) com $g\equiv\E[Y]$, com o lado esquerdo
igual a $\Var(Y)$; 0,2 por usar $\E[r(\X)]=\E[Y]$ e chegar a
$\Var(Y)-\risco(r)=\Var\big(r(\X)\big)$; 0,3 pela interpretação: é o máximo que se
ganha trocando o preditor constante pelo melhor preditor possível, quanto a resposta
esperada varia de um perfil de covariáveis para outro. Escrever a lei da variância total
sem passar pela identidade de (a) vale 0,3 dos 0,5 da conta; a leitura como $R^2$
populacional não é exigida.
\end{criterio}

%==============================================================================
\questao{O que o MQO estima quando a linearidade é falsa}{3,0 pontos}
%==============================================================================
\pontositens{0,7;0,8;0,7;0,8}
% Aula 02 --- o melhor preditor linear (o "oráculo") quando r não é linear.
Circula por aí a ideia de que a regressão linear "só vale se o modelo estiver
correto". Vimos durante o curso que isso é impreciso: o MQO converge para uma função
bem definida mesmo quando $r$ não é linear.

Essa função é o \textbf{melhor preditor linear},
\[
  \bbeta^\ast = \argmin_{\bbeta} \E\big[(Y-\bbeta^\top\X)^2\big],
\]
que existe sempre que $\Sigma=\E[\X\X^\top]$ é bem definida e invertível ---
\emph{sem} supor nada sobre a forma de $r$.
\begin{itens}
  \item \pts{0,9} Derive $\bbeta^\ast=\Sigma^{-1}\bm\alpha$, onde
        $\bm\alpha=\E[Y\X]$. É o mesmo cálculo das equações normais, com
        esperanças no lugar de somas.
  \item \pts{1,0} O estimador de mínimos quadrados pode ser escrito como
        $\widehat\bbeta=\widehat\Sigma^{-1}\widehat{\bm\alpha}$, com
        $\widehat\Sigma=\frac1n\sum_i\X_i\X_i^\top$ e
        $\widehat{\bm\alpha}=\frac1n\sum_i Y_i\X_i$. Olhe para essas duas
        expressões e diga o que elas são, estatisticamente. Em seguida, explique
        por que $\widehat\bbeta\to\bbeta^\ast$ em probabilidade sem que se precise
        supor linearidade em lugar nenhum.
  \item \pts{1,0} Todo teorema garante uma coisa e deixa de garantir outra.
        Enuncie as duas, com cuidado. Em particular: o resultado diz que o risco
        de $\widehat\bbeta$ é pequeno?
  \item \pts{1,1} Um colega ajusta uma regressão linear a dados em que $r$ é
        visivelmente curva e conclui: "o coeficiente deu $2{,}3$, logo aumentar
        $x_1$ em uma unidade aumenta $Y$ em $2{,}3$". Separe o que é leitura
        legítima de $\bbeta^\ast$ do que não é.
\end{itens}

\begin{solucao}
Hipóteses em uso: $\E[Y^2]<\infty$ e $\E\norm{\X}^2<\infty$, para que $\Sigma$ e
$\bm\alpha$ existam (a existência de $\bm\alpha$ sai por Cauchy--Schwarz), e $\Sigma$
invertível. Como $\bm v^\top\Sigma\bm v=\E\big[(\bm v^\top\X)^2\big]\ge0$ para todo
$\bm v$, $\Sigma$ é semidefinida positiva; sendo invertível, é definida positiva.

\smallskip
\textbf{(a)} Expanda o objetivo; as esperanças são números fixos:
\[
  Q(\bbeta)=\E\big[(Y-\bbeta^\top\X)^2\big]
  =\E[Y^2]-2\,\bbeta^\top\underbrace{\E[Y\X]}_{\bm\alpha}
   +\bbeta^\top\underbrace{\E[\X\X^\top]}_{\Sigma}\bbeta .
\]
É uma função quadrática de $\bbeta$, com gradiente $\nabla Q(\bbeta)=-2\bm\alpha+2\Sigma\bbeta$
e Hessiana $2\Sigma$, definida positiva. Logo $Q$ é estritamente convexa, e seu único
ponto crítico é o mínimo global:
\[
  \Sigma\bbeta=\bm\alpha
  \iff
  \bbeta^\ast=\Sigma^{-1}\bm\alpha .
\]
São as equações normais da aula com $\E[\X\X^\top]$ no lugar de
$\frac1n\mathbb{X}^\top\mathbb{X}$ e $\E[Y\X]$ no lugar de $\frac1n\mathbb{X}^\top\bm y$:
o mesmo problema, escrito na população em vez de na amostra.

Vale conhecer uma segunda demonstração, sem derivadas, porque ela rende uma identidade
que usaremos em (c). Some e subtraia $\bbeta^{\ast\top}\X$ dentro do quadrado:
\[
\begin{aligned}
  Q(\bbeta)
  &=\E\Big[\big((Y-\bbeta^{\ast\top}\X)-(\bbeta-\bbeta^\ast)^\top\X\big)^2\Big]\\
  &=Q(\bbeta^\ast)
   -2(\bbeta-\bbeta^\ast)^\top\underbrace{\E\big[\X\,(Y-\X^\top\bbeta^\ast)\big]}_{=\;\bm\alpha-\Sigma\bbeta^\ast\;=\;\bm 0}
   +(\bbeta-\bbeta^\ast)^\top\Sigma\,(\bbeta-\bbeta^\ast),
\end{aligned}
\]
ou seja,
\[
  Q(\bbeta)=Q(\bbeta^\ast)+(\bbeta-\bbeta^\ast)^\top\Sigma\,(\bbeta-\bbeta^\ast).
  \tag{$\dagger$}
\]
O último termo é estritamente positivo para $\bbeta\neq\bbeta^\ast$, o que prova de uma
vez que $\bbeta^\ast$ é o mínimo e que é o único. É o mesmo "some e subtraia" que
mostra que a função de regressão é o melhor preditor, e funciona pela mesma razão. Lá,
o resíduo $Y-r(\X)$ é ortogonal a toda função de $\X$; aqui, $\E\big[\X(Y-\X^\top\bbeta^\ast)\big]=\bm0$ diz
que o resíduo do melhor preditor linear é ortogonal a cada covariável. $r(\X)$ é a
projeção de $Y$ sobre todas as funções de $\X$; $\bbeta^{\ast\top}\X$, a projeção sobre
as lineares.

\smallskip
\textbf{(b)} $\widehat\Sigma$ e $\widehat{\bm\alpha}$ são \textbf{médias amostrais} de
variáveis i.i.d.: a primeira, das matrizes $\X_i\X_i^\top$, cuja esperança é $\Sigma$;
a segunda, dos vetores $Y_i\X_i$, cuja esperança é $\bm\alpha$. É essa observação que
resolve o item.

Pela lei dos grandes números, aplicada coordenada a coordenada (as esperanças existem
pelas hipóteses acima), $\widehat\Sigma\to\Sigma$ e $\widehat{\bm\alpha}\to\bm\alpha$ em
probabilidade. A função $(A,\bm a)\mapsto A^{-1}\bm a$ é contínua no conjunto aberto das
matrizes invertíveis, que contém $\Sigma$; por isso $\widehat\Sigma$ é invertível com
probabilidade tendendo a $1$, e o teorema de Mann--Wald (da aplicação contínua)
transporta a convergência:
\[
  \widehat\bbeta=\widehat\Sigma^{-1}\widehat{\bm\alpha}
  \;\xrightarrow{\ P\ }\;\Sigma^{-1}\bm\alpha=\bbeta^\ast .
\]
Em nenhum passo apareceu a forma de $r$. O alvo $\bbeta^\ast$ é definido por um
problema de minimização que faz sentido para qualquer distribuição com segundos
momentos finitos; o que a linearidade de $r$ acrescentaria seria só a coincidência
$\bbeta^{\ast\top}\x=r(\x)$ --- agradável, mas dispensável para a convergência.

\smallskip
\textbf{(c)} Duas identidades dizem tudo. A primeira vale para qualquer função $g$ com
$\E[g(\X)^2]<\infty$:
\[
  \E\big[(Y-g(\X))^2\big]=\E\big[\Var(Y\mid\X)\big]+\E\big[(r(\X)-g(\X))^2\big],
\]
e sai de somar e subtrair $r(\X)$ dentro do quadrado: o termo cruzado se anula pela
esperança total, porque $\E[Y-r(\X)\mid\X]=0$, e $\E[(Y-r(\X))^2]=\E[\Var(Y\mid\X)]$.
Aplique-a com $g(\x)=\bbeta^{\ast\top}\x$, e depois a segunda, $(\dagger)$, com
$\bbeta=\widehat\bbeta$ (fixado o treino, o risco do ajuste é
$\risco(\widehat g)=Q(\widehat\bbeta)$):
\[
  \risco(\widehat g)=
  \underbrace{\E\big[\Var(Y\mid\X)\big]}_{\text{irredutível}}
  +\underbrace{\E\big[(r(\X)-\bbeta^{\ast\top}\X)^2\big]}_{\text{erro de aproximação}}
  +\underbrace{(\widehat\bbeta-\bbeta^\ast)^\top\Sigma\,(\widehat\bbeta-\bbeta^\ast)}_{\text{erro de estimação}} .
\]
Com a decomposição à vista, as duas coisas se enunciam sem ambiguidade.

\textbf{O que o teorema garante:} a terceira parcela vai a zero, porque
$\widehat\bbeta\to\bbeta^\ast$ e a forma quadrática é contínua. Com $n$ grande, o MQO
chega tão perto quanto se queira do melhor que a classe das funções lineares tem a
oferecer, e $\risco(\widehat g)\to\risco(g_{\bbeta^\ast})$.

\textbf{O que ele não garante:} que esse melhor seja bom. A segunda parcela não depende
de $n$, e o teorema não diz nada sobre ela; se $r$ é muito curva, ela é grande, e
convergir para uma aproximação ruim continua sendo ruim. Então não: o resultado
\emph{não} diz que o risco de $\widehat\bbeta$ é pequeno, só que ele se aproxima do
risco do oráculo. Tampouco diz \emph{com que rapidez}: é um resultado de convergência,
sem taxa.

Na linguagem do balanço entre viés e variância, o teorema cuida do erro de estimação --- a variância, e a
parte do viés que vem de $n$ ser finito --- e deixa intacto o viés de aproximação, que
é o preço fixo de ter escolhido uma classe rígida.

\smallskip
\textbf{(d)} \emph{A leitura legítima.} $2{,}3$ é o coeficiente de $x_1$ no hiperplano
que \emph{melhor aproxima} $r$ em erro quadrático médio, sob a distribuição de $\X$ em
que os dados foram colhidos. É um resumo global, e útil como tal.

\emph{O que não se pode ler nele}, em três sentidos.
\begin{itemize}[leftmargin=1.2cm]
  \item \emph{Não é a inclinação de $r$.} Com $r$ curva, $\partial r/\partial x_1$ muda
        de ponto para ponto, e $2{,}3$ não é o valor dela em lugar nenhum em
        particular. Com várias covariáveis, pode nem ter relação com ela. Tome
        $X_2\sim N(0,1)$, $X_1=X_2^2+\varepsilon$, com $\varepsilon\sim N(0,s^2)$
        independente, e $r(\x)=x_2^2$: $r$ nem depende de $x_1$, e ainda assim,
        regredindo em $(1,x_1,x_2)$,
        \[
          \beta_1^\ast=\frac{\Cov(X_1,X_2^2)}{\Var(X_1)}=\frac{2}{2+s^2}\neq0 .
        \]
        A reta usa $x_1$ como \emph{procurador} da parte curva de $r$ que ela não
        consegue representar em $x_2$. (A conta: $X_2$ não é correlacionada nem com
        $X_2^2$ nem com $X_1$, pois $\E[X_2^3]=0$; então o coeficiente de $x_2$ é zero,
        o de $x_1$ é o da regressão simples, e $\Var(X_2^2)=\E[X_2^4]-1=2$.)
  \item \emph{Depende de onde os dados caíram.} Mesmo com uma covariável só: para
        $r(x)=x^2$ e $X\sim\mathrm{Unif}(0,a)$, a inclinação ótima é
        $\Cov(X,X^2)/\Var(X)=(a^3/12)/(a^2/12)=a$. Colhendo dados em $[0,1]$ o colega
        veria $1$; em $[0,2]$, veria $2$ --- com $r$ exatamente a mesma. Um número que
        muda com a faixa amostrada não descreve um mecanismo.
  \item \emph{"Aumentar $x_1$ aumenta $Y$" é uma frase causal}, sobre intervir no
        sistema. Nenhuma regressão a autoriza sozinha, nem com $r$ linear: comparar
        indivíduos que \emph{diferem} em $x_1$ não é o mesmo que \emph{mudar} o $x_1$
        de alguém.
\end{itemize}
\end{solucao}


\begin{criterio}
\textbf{(a)} 0,3 pela expansão
$Q(\bbeta)=\E[Y^2]-2\,\bbeta^\top\bm\alpha+\bbeta^\top\Sigma\bbeta$; 0,2 pelo gradiente
igualado a zero, $\Sigma\bbeta=\bm\alpha$; 0,2 por justificar que o ponto crítico é o
mínimo, pela Hessiana $2\Sigma$ definida positiva ou pela identidade $(\dagger)$ da
solução. Aceitar integralmente a demonstração que só usa $(\dagger)$.

\textbf{(b)} 0,3 por reconhecer $\widehat\Sigma$ e $\widehat{\bm\alpha}$ como médias
amostrais de variáveis i.i.d., com esperanças $\Sigma$ e $\bm\alpha$; 0,2 pela lei dos
grandes números; 0,2 pela passagem a $\widehat\Sigma^{-1}\widehat{\bm\alpha}$, pela
continuidade da inversão no conjunto das matrizes invertíveis; 0,1 por observar que a
forma de $r$ não aparece em nenhum passo. Não é exigido o nome do teorema da aplicação
contínua.

\textbf{(c)} 0,3 pelo que o teorema garante: o erro de estimação vai a zero, e o risco do
ajuste se aproxima do risco do melhor preditor linear; 0,3 pelo que ele não garante:
que esse risco seja pequeno, porque o erro de aproximação
$\E[(r(\X)-\bbeta^{\ast\top}\X)^2]$ não depende de $n$; 0,1 pela resposta à pergunta
final, que é não. A decomposição em três parcelas é a forma mais clara de responder, mas
não é exigida, e a ausência de taxa também não.

\textbf{(d)} 0,3 pela leitura legítima: $2{,}3$ é o coeficiente de $x_1$ no hiperplano que
melhor aproxima $r$ em erro quadrático médio, sob a distribuição de $\X$ em que os dados
foram colhidos. 0,5 pelo que não se pode ler, em pelo menos dois dos três sentidos da
solução, 0,25 cada: não é a inclinação de $r$; depende da região em que os dados
caíram; ``aumentar $x_1$ aumenta $Y$'' é uma afirmação causal. Os exemplos numéricos da
solução não são exigidos.
\end{criterio}

%==============================================================================
\questao{}{3,0 pontos}
%==============================================================================
\pontositens{0,7;0,8;0,9;0,6}
% Aula 07 --- uma covariável duplicada: o Bayes ingênuo conta a mesma evidência duas vezes.
Lucas é analista de crédito e quer prever a inadimplência, $Y=1$, a partir da renda
mensal $X_1$ do cliente, em milhares de reais. Na população de clientes,
$X_1\mid Y=1\sim N(3,1)$ e $X_1\mid Y=0\sim N(5,1)$. Ao juntar duas bases, Lucas fica,
sem perceber, com a renda duas vezes: em $X_1$, em milhares de reais, e em $X_2=1000\,X_1$,
em reais. Ele ajusta um Bayes ingênuo gaussiano às duas colunas. Para separar o erro do
modelo do erro de estimação, suponha que a amostra é tão grande que o ingênuo usa as
distribuições verdadeiras de $X_1$ e de $X_2$ em cada classe. Seja
\[
  L(x_1)=\log\frac{f_1(x_1)}{f_0(x_1)}
\]
a razão de log-verossimilhanças da renda, em que $f_c$ é a densidade de $X_1$ dado $Y=c$
e $\log$ é o logaritmo natural, e sejam $\pi_c=\Prob(Y=c)$ e $\sigma(u)=e^u/(1+e^u)$.
\begin{itens}
  \item \pts{1,0} Mostre que $L(x_1)=8-2x_1$. Mostre que, num cliente com renda $x_1$, o
        \emph{log-odds} verdadeiro é $\log(\pi_1/\pi_0)+L(x_1)$ e o calculado pelo
        ingênuo é $\log(\pi_1/\pi_0)+2L(x_1)$.
  \item \pts{1,0} Suponha $\pi_0=\pi_1=1/2$ e seja $\eta=\Prob(Y=1\mid X_1=x_1)$. Mostre
        que a probabilidade calculada pelo ingênuo é
        $\widehat P=\eta^2/\big(\eta^2+(1-\eta)^2\big)$. Conclua que ela ordena os
        clientes como $\eta$, que classifica cada um como o classificador de Bayes, e que
        $\abs{\widehat P-1/2}>\abs{\eta-1/2}$ sempre que $\eta\notin\{0,1/2,1\}$.
        Calcule $\widehat P$ para $\eta=0{,}6$ e para $\eta=0{,}8$.
  \item \pts{1,2} Na carteira de Lucas, porém, só $20\%$ dos clientes são inadimplentes:
        $\pi_1=0{,}2$. Determine as rendas que o classificador de Bayes classifica como
        inadimplentes e as que o ingênuo classifica assim, e mostre que as duas decisões
        diferem num intervalo de rendas. Calcule a probabilidade de um cliente cair nesse
        intervalo e diga qual das duas decisões erra mais nele. \emph{Dados:}
        $\log2\approx0{,}6931$, $\Phi(-1{,}3466)\approx0{,}0891$,
        $\Phi(-1{,}6931)\approx0{,}0452$, $\Phi(0{,}6534)\approx0{,}7433$ e
        $\Phi(0{,}3069)\approx0{,}6205$, em que $\Phi$ é a função de distribuição da
        normal padrão.
  \item \pts{0,8} Se a renda aparecesse em $m$ colunas, todas proporcionais a $X_1$, o que
        aconteceria com o \emph{log-odds} do ingênuo? E com a sua decisão quando
        $m\to\infty$? Resuma o que o exercício mostra sobre o Bayes ingênuo com
        covariáveis redundantes: o que ele preserva e o que ele estraga.
\end{itens}

\begin{solucao}
Em todo o exercício, $f_1$ e $f_0$ são as densidades de $N(3,1)$ e de $N(5,1)$, e
$\pi_1=1-\pi_0\in(0,1)$. Como pede o enunciado, o ingênuo usa as distribuições verdadeiras de
cada coluna em cada classe; assim, toda diferença entre ele e o classificador de Bayes vem da
hipótese de independência condicional, e não da estimação. Os dois declaram inadimplente
quando o seu \emph{log-odds} é positivo, isto é, quando a probabilidade que calculam passa de
$1/2$.

\smallskip
\textbf{(a)} As duas normais têm variância $1$, e a constante $1/\sqrt{2\pi}$ se cancela no
quociente:
\[
  L(x_1)=-\tfrac12(x_1-3)^2+\tfrac12(x_1-5)^2
        =\tfrac12\big[(x_1-5)^2-(x_1-3)^2\big]
        =\tfrac12\,(-2)\,(2x_1-8)=8-2x_1 ,
\]
em que se usou $a^2-b^2=(a-b)(a+b)$, com $a=x_1-5$ e $b=x_1-3$, de modo que $a-b=-2$ e
$a+b=2x_1-8$.

\emph{O log-odds verdadeiro.} Como $X_2=1000\,X_1$, o par $(X_1,X_2)$ é uma função de $X_1$, e
$X_1=X_2/1000$ é uma função do par: observar o par é o mesmo que observar $X_1$, e a
distribuição de $Y$ dado o par é a distribuição de $Y$ dado $X_1$. Pelo Teorema de Bayes, num
cliente com renda $x_1$,
\[
  \log\frac{\Prob(Y=1\mid X_1=x_1)}{\Prob(Y=0\mid X_1=x_1)}
  =\log\frac{\pi_1f_1(x_1)}{\pi_0f_0(x_1)}=\log\frac{\pi_1}{\pi_0}+L(x_1).
\]

\emph{O log-odds do ingênuo.} O ingênuo supõe que, dada a classe, as colunas são
independentes, e usa como densidade do par o produto das densidades marginais,
$f_c(x_1)\,h_c(x_2)$, em que $h_c$ é a densidade de $X_2$ dado $Y=c$. Como
$\Prob(X_2\le t\mid Y=c)=\Prob(X_1\le t/1000\mid Y=c)$, derivar em $t$ dá
\[
  h_c(t)=\frac1{1000}\,f_c\Big(\frac t{1000}\Big),
\]
que é a fórmula da mudança de variável: $X_2\mid Y=c\sim N(1000\,\mu_c,\,1000^2)$. No ponto
observado, $x_2=1000\,x_1$, o fator $1/1000$ aparece no numerador e no denominador, e
$h_1(1000\,x_1)/h_0(1000\,x_1)=f_1(x_1)/f_0(x_1)$. O \emph{log-odds} do ingênuo é, então,
\[
  \log\frac{\pi_1f_1(x_1)\,h_1(1000\,x_1)}{\pi_0f_0(x_1)\,h_0(1000\,x_1)}
  =\log\frac{\pi_1}{\pi_0}+2L(x_1).
\]
A mudança de unidade não altera a razão de verossimilhanças. O que o ingênuo faz de errado é
tratar $X_2$ como uma segunda evidência, independente da primeira dada a classe, quando ela é a
mesma informação. A hipótese de independência condicional falha aqui no grau máximo, porque
$X_2$ é uma função de $X_1$.

\smallskip
\textbf{(b)} Com $\pi_0=\pi_1=1/2$, $\log(\pi_1/\pi_0)=0$. Escreva $u=L(x_1)$: pelo item (a),
$\eta=\sigma(u)$, e a probabilidade do ingênuo é $\widehat P=\sigma(2u)$. De
$\eta=e^u/(1+e^u)$ vem $e^u=\eta/(1-\eta)$, e
\[
  \widehat P=\frac{e^{2u}}{1+e^{2u}}
  =\frac{\eta^2/(1-\eta)^2}{1+\eta^2/(1-\eta)^2}
  =\frac{\eta^2}{\eta^2+(1-\eta)^2}.
\]

\emph{A ordenação.} Com $\operatorname{logit}\eta=\log[\eta/(1-\eta)]$, vale
$\widehat P=\sigma\big(2\operatorname{logit}\eta\big)$, composição de três funções
estritamente crescentes: o $\operatorname{logit}$, a multiplicação por $2$ e $\sigma$. Logo
$\widehat P$ é uma função estritamente crescente de $\eta$; se um cliente tem $\eta$ maior que
outro, tem também $\widehat P$ maior, e os dois ordenam os clientes do mesmo modo.

\emph{A decisão.} Como $\sigma(v)>1/2$ se e somente se $v>0$,
$\widehat P>1/2\iff 2u>0\iff u>0\iff\eta>1/2$: o ingênuo declara inadimplente exatamente quem
o classificador de Bayes declara.

\emph{O exagero.} Subtraindo,
\[
  \widehat P-\eta
  =\frac{\eta^2-\eta\big(\eta^2+(1-\eta)^2\big)}{\eta^2+(1-\eta)^2}
  =\frac{\eta\,(1-\eta)\,(2\eta-1)}{\eta^2+(1-\eta)^2},
\]
porque $\eta-\eta^2-(1-\eta)^2=-2\eta^2+3\eta-1=(1-\eta)(2\eta-1)$. Para $0<\eta<1$, o fator
$\eta(1-\eta)$ e o denominador são positivos, e $\widehat P-\eta$ tem o sinal de $2\eta-1$. Se
$\eta>1/2$, então $\widehat P>\eta>1/2$; se $\eta<1/2$, então $\widehat P<\eta<1/2$. Nos dois
casos, $\abs{\widehat P-1/2}>\abs{\eta-1/2}$: a probabilidade do ingênuo fica do mesmo lado de
$1/2$ que a verdadeira, e mais longe. Com $\eta=0{,}6$,
$\widehat P=0{,}36/0{,}52=9/13\approx0{,}692$; com $\eta=0{,}8$,
$\widehat P=0{,}64/0{,}68=16/17\approx0{,}941$.

Com prioris iguais, portanto, a coluna repetida não muda nenhuma decisão nem a ordem dos
clientes, mas muda a probabilidade informada: um cliente com $60\%$ de chance de
inadimplência aparece com $69\%$, e um com $80\%$, com $94\%$.

\smallskip
\textbf{(c)} Com $\pi_1=0{,}2$, $\log(\pi_1/\pi_0)=\log(1/4)=-2\log2$. O classificador de Bayes
declara inadimplente quando o \emph{log-odds} verdadeiro é positivo,
\[
  -2\log2+8-2x_1>0\iff x_1<4-\log2=:a\approx3{,}3069 ,
\]
e o ingênuo, quando o seu é positivo,
\[
  -2\log2+2\,(8-2x_1)>0\iff x_1<4-\tfrac12\log2=:b\approx3{,}6534 .
\]
Como $a<b$, os dois declaram inadimplente quem tem $x_1<a$, nenhum declara quem tem $x_1\ge b$,
e as decisões diferem exatamente em $I=[a,b)$, as rendas entre cerca de R\$~3\,307 e
R\$~3\,653: ali o ingênuo declara inadimplente quem o classificador de Bayes declara
adimplente. Com prioris iguais isso não acontecia, porque multiplicar o \emph{log-odds} por $2$
não muda o seu sinal. Com prioris diferentes, o termo $\log(\pi_1/\pi_0)$ não é dobrado, só a
evidência é, e o ponto em que a soma troca de sinal se desloca de $a$ para $b$.

Como $X_1\mid Y=c\sim N(\mu_c,1)$, com $\mu_0=5$ e $\mu_1=3$, vale
$\Prob(X_1\in I\mid Y=c)=\Phi(b-\mu_c)-\Phi(a-\mu_c)$. Com os dados do enunciado,
\begin{align*}
  \Prob(X_1\in I\mid Y=0)&=\Phi(-1{,}3466)-\Phi(-1{,}6931)\approx0{,}0891-0{,}0452=0{,}0439,\\
  \Prob(X_1\in I\mid Y=1)&=\Phi(0{,}6534)-\Phi(0{,}3069)\approx0{,}7433-0{,}6205=0{,}1228,
\end{align*}
e, pela lei da probabilidade total,
\[
  \Prob(X_1\in I)\approx0{,}8\cdot0{,}0439+0{,}2\cdot0{,}1228\approx0{,}0351+0{,}0246\approx0{,}060 :
\]
a decisão muda para cerca de $6\%$ dos clientes.

Em $I$, o ingênuo declara inadimplente e erra nos adimplentes, com probabilidade
$\Prob(Y=0,\,X_1\in I)\approx0{,}0351$; o classificador de Bayes declara adimplente e erra nos
inadimplentes, com probabilidade $\Prob(Y=1,\,X_1\in I)\approx0{,}0246$. Fora de $I$ as duas
decisões coincidem, e os erros também. O ingênuo erra mais, e a diferença entre as duas
probabilidades de erro é
\[
  \Prob(Y=0,\,X_1\in I)-\Prob(Y=1,\,X_1\in I)\approx0{,}0351-0{,}0246\approx0{,}011 .
\]
Ela tinha de ser positiva, porque nenhum classificador tem probabilidade de erro menor que a do
classificador de Bayes. Para ter a escala, uma regra que declara inadimplente quem tem $x_1<t$
erra com probabilidade $0{,}8\,\Phi(t-5)+0{,}2\,\big[1-\Phi(t-3)\big]$, e, com os mesmos dados,
\begin{align*}
  \text{classificador de Bayes } (t=a)&:\quad 0{,}8\cdot0{,}0452+0{,}2\cdot(1-0{,}6205)\approx0{,}112,\\
  \text{ingênuo } (t=b)&:\quad 0{,}8\cdot0{,}0891+0{,}2\cdot(1-0{,}7433)\approx0{,}123 .
\end{align*}

\smallskip
\textbf{(d)} Sejam as colunas $c_1X_1,\dots,c_mX_1$, com constantes $c_j\ne0$ (no enunciado,
$c_1=1$ e $c_2=1000$). Pela mudança de variável, a densidade de $c_jX_1$ dado $Y=c$ é
$t\mapsto\abs{c_j}^{-1}f_c(t/c_j)$, e o argumento do item (a) se repete: o fator
$\abs{c_j}^{-1}$ se cancela, e cada coluna contribui com $L(x_1)$ para o \emph{log-odds} do
ingênuo, que passa a
\[
  \log\frac{\pi_1}{\pi_0}+m\,L(x_1).
\]
Fixada a renda $x_1$, com $L(x_1)\ne0$, o termo $m\,L(x_1)$ vai a $+\infty$ ou a $-\infty$
quando $m\to\infty$, conforme o sinal de $L(x_1)$, e domina a constante $\log(\pi_1/\pi_0)$; a
probabilidade calculada vai a $1$ se $x_1<4$ e a $0$ se $x_1>4$. A decisão pode ser escrita
para cada $m$: com $\pi_1=0{,}2$, o ingênuo declara inadimplente quando
\[
  -2\log2+m\,(8-2x_1)>0\iff x_1<4-\frac{\log2}{m},
\]
corte que vale $a$ para $m=1$, $b$ para $m=2$, e cresce até $4$ quando $m\to\infty$. No limite,
a decisão é $x_1<4$, a do classificador de Bayes com prioris iguais: a evidência é contada $m$
vezes e a priori, uma vez só, e a informação de que só $20\%$ dos clientes são inadimplentes
deixa de pesar.

O exercício mostra, então, o que o Bayes ingênuo faz com covariáveis redundantes. Preserva a
ordenação: o seu \emph{log-odds}, $\log(\pi_1/\pi_0)+m\,L(x_1)$, e o verdadeiro,
$\log(\pi_1/\pi_0)+L(x_1)$, são funções afins estritamente decrescentes de $x_1$, e ordenam os
clientes do mesmo modo. Estraga as probabilidades, que saem empurradas para $0$ e $1$, e, quando
as classes têm tamanhos diferentes, estraga também as decisões numa faixa de rendas, que cresce
com $m$.

Colunas proporcionais são o caso extremo, e a conclusão não se transfere sem cuidado para
covariáveis apenas correlacionadas. Para ver o que muda, suponha prioris iguais e
$\X\mid Y=c\sim N(\bm\mu_c,\bm\Sigma)$ em $\R^2$, com $\bm\mu_0=-\bm\mu_1$, variâncias $1$ e
correlação $\rho$, $0<\abs\rho<1$. O \emph{log-odds} verdadeiro é linear, com coeficiente
$\bm\Sigma^{-1}(\bm\mu_1-\bm\mu_0)$, como no LDA, e o do ingênuo, que usa só as variâncias, tem
coeficiente $\bm\mu_1-\bm\mu_0$; com médias opostas, os dois interceptos se anulam. Se as médias
diferem de $d$ nas duas coordenadas, os coeficientes são $\frac d{1+\rho}(1,1)$ e $d\,(1,1)$: o
\emph{log-odds} do ingênuo é $1+\rho$ vezes o verdadeiro, o que exagera as probabilidades quando
$\rho>0$ e as atenua quando $\rho<0$. Se as médias diferem de $d\ne0$ só na primeira coordenada,
os coeficientes são $\frac d{1-\rho^2}(1,-\rho)$ e $d\,(1,0)$. O verdadeiro usa $X_2$, que tem a
mesma distribuição nas duas classes mas ajuda a classificar pela correlação com $X_1$, e o
ingênuo a ignora: duas observações com o mesmo $x_1$ e valores diferentes de $x_2$ empatam no
ingênuo, mas não no \emph{log-odds} verdadeiro, e as duas ordenações diferem. Com colunas
repetidas, a conduta é retirá-las antes do ajuste; com colunas muito correlacionadas,
retirá-las ou usar um modelo que leve a correlação em conta.
\end{solucao}


\begin{criterio}
\textbf{(a)} 0,2 por $L(x_1)=8-2x_1$; 0,2 pelo \emph{log-odds} verdadeiro, com a
justificativa de que observar o par é observar $X_1$; 0,3 pelo do ingênuo: 0,1 pelo
produto das densidades marginais e 0,2 pela densidade de $X_2=1000\,X_1$, pela mudança de
variável, cujo fator $1/1000$ se cancela. Afirmar que a segunda coluna repete a
evidência, sem a conta, vale 0,1 desses 0,3.

\textbf{(b)} 0,2 pela fórmula de $\widehat P$, a partir de $e^u=\eta/(1-\eta)$; 0,2 pela
ordenação ($\widehat P$ é função estritamente crescente de $\eta$) e pela decisão
($\widehat P>1/2\iff\eta>1/2$); 0,3 por $\abs{\widehat P-1/2}>\abs{\eta-1/2}$, pela
identidade $\widehat P-\eta=\eta(1-\eta)(2\eta-1)/\big(\eta^2+(1-\eta)^2\big)$ ou pelo
caso $\eta>1/2$ com a simetria; 0,1 pelos dois valores, $0{,}692$ e $0{,}941$.

\textbf{(c)} 0,3 pelos dois cortes, $x_1<4-\log2\approx3{,}3069$ e
$x_1<4-\tfrac12\log2\approx3{,}6534$, e pelo intervalo $I$ em que as decisões diferem;
0,3 por $\Prob(X_1\in I)\approx0{,}060$, com as duas probabilidades condicionais e a lei
da probabilidade total; 0,3 por quem erra mais em $I$: o ingênuo, com
$\Prob(Y=0,\,X_1\in I)\approx0{,}035$, contra $\Prob(Y=1,\,X_1\in I)\approx0{,}025$ do
classificador de Bayes. Explicar que, com prioris diferentes, o termo $\log(\pi_1/\pi_0)$
não é dobrado é bem-vindo, mas não é exigido.

\textbf{(d)} 0,2 pelo \emph{log-odds} $\log(\pi_1/\pi_0)+m\,L(x_1)$; 0,2 pelo limite:
probabilidades indo a $0$ ou a $1$ e decisão $x_1<4$, a de prioris iguais; 0,2 pelo
resumo: o ingênuo preserva a ordenação, estraga as probabilidades e, com classes de
tamanhos diferentes, algumas decisões. A discussão de covariáveis apenas correlacionadas,
que a solução traz, não é exigida.
\end{criterio}

%==============================================================================
\questao{}{3,0 pontos}
%==============================================================================
\pontositens{0,9;1,2;0,9}
% Aula 09 --- a margem máxima depende só dos vetores de suporte, e o LOOCV fica limitado por |S|/n.
Tiago trabalha no controle de qualidade de uma fábrica e classifica peças em boas ($y=+1$) e
defeituosas ($y=-1$) a partir de $p$ medidas. As suas $n$ observações de treino
$(\x_i,y_i)$ são linearmente separáveis, e cada classe tem pelo menos duas observações. O
classificador de margem máxima é dado pela solução do problema
\[
  \max_{\beta_0,\,\bbeta,\,M}\ M
  \quad\text{sujeito a}\quad
  \sum_{j=1}^p\beta_j^2=1
  \quad\text{e}\quad
  y_i\,(\beta_0+\bbeta^\top\x_i)\ge M,\ \ i=1,\dots,n .
\]
Defina $f(\x)=\beta_0+\bbeta^\top\x$. Os vetores de suporte são as observações com
$y_if(\x_i)=M$, e $\mathcal S$ é o conjunto dos seus índices. Admita que, para observações
separáveis com as duas classes presentes, o problema tem solução, e seja $M^*>0$ o valor
ótimo. Durante as aulas afirmamos, sem demonstrar, que essa solução depende apenas dos vetores de
suporte. Este exercício pretende elucidar esse resultado.
\begin{itens}
  \item \pts{1,2} Mostre que a solução é única. (Se houvesse duas soluções com vetores
        $\bbeta$ diferentes, compare-as com o ponto médio entre elas. Se tivessem o mesmo
        $\bbeta$, olhe, como funções de $\beta_0$, o menor valor de $y_if(\x_i)$ em cada
        classe.)
  \item \pts{1,6} Seja $i_0\notin\mathcal S$. Mostre que o classificador de margem máxima
        das $n-1$ observações sem a observação $i_0$ é o mesmo das $n$ observações. (Se o
        valor ótimo sem $i_0$ fosse maior que $M^*$, combine a solução nova com a antiga, com
        peso pequeno na nova.) Mostre também que deslocar $\x_{i_0}$ para qualquer ponto
        $\x'$ com $y_{i_0}f(\x')>M^*$, em que $f$ é a função da solução, não muda o
        classificador.
  \item \pts{1,2} A estimativa do risco por validação cruzada com uma observação deixada de
        fora (LOOCV) é a proporção das observações $i$ que o classificador ajustado sem $i$
        classifica errado. Mostre que, para o classificador de margem máxima, ela é no
        máximo $\abs{\mathcal S}/n$.
\end{itens}

\begin{solucao}
Para $(b,\bm w)\in\R\times\R^p$, escreva $f_{b,\bm w}(\x)=b+\bm w^\top\x$. Fixados $(b,\bm w)$,
as restrições permitem qualquer $M\le\min_iy_if_{b,\bm w}(\x_i)$, e o problema maximiza, entre
os $(b,\bm w)$ com $\norm{\bm w}=1$, a margem
\[
  m(b,\bm w)=\min_{1\le i\le n}\,y_i\,f_{b,\bm w}(\x_i);
\]
$(\beta_0,\bbeta,M)$ é solução se e somente se $\norm{\bbeta}=1$ e
$M=m(\beta_0,\bbeta)=M^*$. Dois fatos simples serão usados várias vezes.
\begin{itemize}[leftmargin=1.2cm]
  \item[(F1)] Se $y_if_{b,\bm w}(\x_i)>0$ para todo $i$ e as duas classes estão presentes,
        então $\bm w\ne\bm0$. De fato, se $\bm w=\bm0$, então $y_ib>0$ para todo $i$, e $b$
        teria de ser positivo, por uma observação com $y_i=+1$, e negativo, por uma com
        $y_i=-1$.
  \item[(F2)] Se $0<\norm{\bm w}\le1$ e $y_if_{b,\bm w}(\x_i)\ge\mu>0$ para todo $i$, então
        o par $(b,\bm w)/\norm{\bm w}$ satisfaz a restrição de norma e tem margem pelo menos
        $\mu/\norm{\bm w}\ge\mu$, porque dividir números positivos por $\norm{\bm w}\le1$ não
        os diminui; se $\norm{\bm w}<1$, a margem é estritamente maior que $\mu$.
\end{itemize}

\smallskip
\textbf{(a)} \emph{Unicidade de $\bbeta$.} Sejam $(\beta_0,\bbeta,M^*)$ e $(\beta_0',\bbeta',M^*)$
duas soluções, e suponha $\bbeta\ne\bbeta'$. Seja $(\bar b,\bar{\bm w})$ o ponto médio de
$(\beta_0,\bbeta)$ e $(\beta_0',\bbeta')$. Como $f_{\bar b,\bar{\bm w}}$ é a média de
$f_{\beta_0,\bbeta}$ e $f_{\beta_0',\bbeta'}$, para todo $i$
\[
  y_i\,f_{\bar b,\bar{\bm w}}(\x_i)
  =\tfrac12\big[y_i\,f_{\beta_0,\bbeta}(\x_i)+y_i\,f_{\beta_0',\bbeta'}(\x_i)\big]\ \ge\ M^*,
\]
porque as duas parcelas são pelo menos $M^*$. Além disso,
$\norm{\bar{\bm w}}^2=\tfrac14\big(\norm{\bbeta}^2+2\,\bbeta^\top\bbeta'+\norm{\bbeta'}^2\big)
=\tfrac12\big(1+\bbeta^\top\bbeta'\big)$. Pela desigualdade de Cauchy--Schwarz,
$\bbeta^\top\bbeta'\le\norm{\bbeta}\,\norm{\bbeta'}=1$, com igualdade somente se
$\bbeta'=c\,\bbeta$ para algum $c\ge0$, o que, com as duas normas iguais a $1$, daria
$\bbeta'=\bbeta$. Logo $\bbeta^\top\bbeta'<1$ e $\norm{\bar{\bm w}}<1$. Por (F1),
$\bar{\bm w}\ne\bm0$, e, por (F2), $(\bar b,\bar{\bm w})/\norm{\bar{\bm w}}$ é viável com margem
maior que $M^*$, o que contradiz a otimalidade de $M^*$. Portanto $\bbeta'=\bbeta$.

\emph{Unicidade de $\beta_0$.} Com $\bbeta$ fixado, separando as observações pela classe,
\begin{align*}
  m(b,\bbeta)
  &=\min\Big\{\min_{i:\,y_i=+1}\big(b+\bbeta^\top\x_i\big),\
              \min_{i:\,y_i=-1}\big(-b-\bbeta^\top\x_i\big)\Big\}\\
  &=\min\{\,b+a_+,\ a_--b\,\},
  \qquad
  a_+=\min_{i:\,y_i=+1}\bbeta^\top\x_i,\quad a_-=-\max_{i:\,y_i=-1}\bbeta^\top\x_i .
\end{align*}
A função $b\mapsto b+a_+$ é estritamente crescente, e $b\mapsto a_--b$, estritamente
decrescente; elas se igualam em $b^*=(a_--a_+)/2$. Para $b<b^*$, o mínimo das duas é
$b+a_+$, que cresce, e para $b>b^*$ é $a_--b$, que decresce: o único ponto de máximo de
$b\mapsto m(b,\bbeta)$ é $b^*$. Logo $\beta_0=\beta_0'=b^*$, e a solução é única.

A fórmula de $b^*$ tem uma leitura geométrica. Na direção $\bbeta$, o hiperplano
$\{\x:\bbeta^\top\x=-b^*\}$ fica exatamente a meio caminho entre a maior projeção de uma
observação da classe $-1$ e a menor projeção de uma da classe $+1$, e a margem é a metade da
distância entre essas duas projeções, $M^*=(a_++a_-)/2$.

\smallskip
\textbf{(b)} Seja $(\beta_0^*,\bbeta^*,M^*)$ a solução com as $n$ observações e
$f^*=f_{\beta_0^*,\bbeta^*}$. Como $i_0\notin\mathcal S$ e as restrições valem,
$\delta=y_{i_0}f^*(\x_{i_0})-M^*>0$. Sem $i_0$, as observações continuam separáveis, por $f^*$,
e com as duas classes, porque cada classe tinha pelo menos duas; pela hipótese do enunciado,
esse problema tem solução, com valor ótimo $M'$. O ponto $(\beta_0^*,\bbeta^*,M^*)$ satisfaz
todas as restrições sem $i_0$, e por isso $M'\ge M^*$.

Suponha $M'>M^*$, atingido por $(\beta_0',\bbeta')$, com $f'=f_{\beta_0',\bbeta'}$. Para
$t\in(0,1]$, seja $(b_t,\bm w_t)=(1-t)(\beta_0^*,\bbeta^*)+t\,(\beta_0',\bbeta')$, de modo que
$f_{b_t,\bm w_t}=(1-t)f^*+t\,f'$. Para $i\ne i_0$,
\[
  y_if_{b_t,\bm w_t}(\x_i)=(1-t)\,y_if^*(\x_i)+t\,y_if'(\x_i)\ \ge\ (1-t)M^*+tM'\ >\ M^* .
\]
Para $i_0$, com $c=\abs{y_{i_0}f'(\x_{i_0})-y_{i_0}f^*(\x_{i_0})}$,
\[
  y_{i_0}f_{b_t,\bm w_t}(\x_{i_0})
  =y_{i_0}f^*(\x_{i_0})+t\big[y_{i_0}f'(\x_{i_0})-y_{i_0}f^*(\x_{i_0})\big]
  \ \ge\ M^*+\delta-t\,c ,
\]
que é maior que $M^*$ sempre que $t\,c<\delta$. Fixe $t\in(0,1]$ com $t\,c<\delta$ (se $c=0$,
qualquer $t$ serve). Então os $n$ números $y_if_{b_t,\bm w_t}(\x_i)$ são todos maiores que
$M^*$, e o menor deles, $\mu$, também é. Pela desigualdade triangular,
$\norm{\bm w_t}\le(1-t)\norm{\bbeta^*}+t\norm{\bbeta'}=1$, e $\bm w_t\ne\bm0$ por (F1). Por
(F2), $(b_t,\bm w_t)/\norm{\bm w_t}$ é viável no problema com as $n$ observações, com margem pelo
menos $\mu>M^*$, o que contradiz a otimalidade de $M^*$. Logo $M'=M^*$: a solução antiga é
viável sem $i_0$ e atinge o valor ótimo, e, pelo item (a), aplicado às $n-1$ observações, é a
única solução sem $i_0$. O classificador não muda.

Para o deslocamento, considere o problema com as $n-1$ observações sem $i_0$ e mais a
observação $(\x',y_{i_0})$, em que $y_{i_0}f^*(\x')>M^*$. Todo ponto viável desse problema
satisfaz as restrições do problema sem $i_0$, e por isso o seu valor ótimo é no máximo
$M'=M^*$. A solução antiga satisfaz as restrições das $n-1$ observações e a da nova, porque
$y_{i_0}f^*(\x')>M^*$, e atinge $M^*$; é, portanto, uma solução. As observações desse problema
são separáveis, por $f^*$, e têm as duas classes, e, pelo item (a), a solução é única: o
classificador também não muda.

As observações do lado certo e fora da faixa fechada $\{\x:\abs{f^*(\x)}\le M^*\}$, isto é,
com $y_if^*(\x_i)>M^*$, não influenciam a solução: podem ser retiradas ou deslocadas, desde que
continuem do lado certo e fora dessa faixa. Só os vetores de suporte determinam o classificador
de margem máxima.

\smallskip
\textbf{(c)} Para cada $i$, as $n-1$ observações sem $i$ são separáveis e têm as duas classes,
porque cada classe tinha pelo menos duas, e o classificador de margem máxima delas, $g_{-i}$,
está bem definido. Se $i\notin\mathcal S$, pelo item (b), $g_{-i}$ é o classificador $g$ das $n$
observações, que classifica $\x_i$ corretamente, porque $y_if^*(\x_i)>M^*>0$. Assim,
$g_{-i}(\x_i)\ne y_i$ só pode acontecer para $i\in\mathcal S$, e
\[
  \frac1n\sum_{i=1}^n\1{g_{-i}(\x_i)\ne y_i}
  =\frac1n\sum_{i\in\mathcal S}\1{g_{-i}(\x_i)\ne y_i}\ \le\ \frac{\abs{\mathcal S}}{n},
\]
porque cada indicadora vale no máximo $1$.

A cota sai de um único ajuste, o das $n$ observações: das $n$ retiradas que a validação cruzada
faria, só as dos vetores de suporte podem alterar o classificador e, com isso, produzir um erro.
Ela é útil quando há poucos vetores de suporte diante de $n$, e pode ser folgada, porque retirar
um vetor de suporte não obriga o novo classificador a errar nele.
\end{solucao}


\begin{criterio}
\textbf{(a)} 0,5 pela unicidade de $\bbeta$: 0,2 pelo ponto médio, viável com margem pelo
menos $M^*$; 0,2 por $\norm{\bar{\bm w}}<1$, pela desigualdade de Cauchy--Schwarz, com
igualdade só se $\bbeta'=\bbeta$; 0,1 pela normalização, que daria margem maior que
$M^*$, com $\bar{\bm w}\ne\bm0$ porque as duas classes estão presentes. 0,4 pela
unicidade de $\beta_0$ com $\bbeta$ fixado, pelo argumento das duas funções monótonas da
solução ou pelo do ponto médio.

\textbf{(b)} 0,2 por $M'\ge M^*$, porque a solução antiga é viável sem $i_0$; 0,6 pelo
absurdo de $M'>M^*$: 0,2 pela combinação convexa e pela desigualdade nas observações
$i\ne i_0$, 0,2 pela escolha de $t$ pequeno, que mantém $i_0$ acima de $M^*$ porque
$\delta>0$, e 0,2 pela normalização ($\norm{\bm w_t}\le1$, pela desigualdade
triangular) e pela contradição; 0,2 pela conclusão, com a unicidade do item (a); 0,2 pelo
deslocamento de $\x_{i_0}$. Seguir a dica sem justificar por que $t$ pequeno funciona
vale 0,1 dos 0,2 correspondentes.

\textbf{(c)} 0,2 por notar que, sem cada observação, os dados continuam separáveis e
com as duas classes; 0,4 pelo uso do item (b): para $i\notin\mathcal S$, o
classificador sem $i$ é o mesmo e acerta $\x_i$, porque $y_if^*(\x_i)>M^*>0$; 0,3 pela
cota, com a soma das indicadoras restrita a $\mathcal S$. A observação de que a cota sai
de um único ajuste é bem-vinda, mas não é exigida.
\end{criterio}

\vfill
\begin{center}\small\itshape Boa prova!\end{center}

\end{document}
